\chapter{Introduction}
\label{ch:introduction}

The foreign exchange (FX) market is the largest and most liquid financial market in the world, with a daily turnover of \$9.6 trillion as of April 2025 \cite{bis2025}. At the heart of this market are dealers, i.e.\ financial institutions that continuously provide liquidity to clients by quoting prices at which they are willing to buy and sell currencies. In doing so, dealers accumulate inventory positions across a wide range of currency pairs, exposing themselves to \emph{inventory risk}: the risk that exchange rate movements adversely affect the value of their holdings before they can be unwound. Managing this risk effectively, while still capturing the spread between bid and ask prices, is the central challenge of FX market making.

Modern FX markets are overwhelmingly electronic, highly fragmented across multiple trading venues, and characterized by diverse client bases and complex order flow dynamics \cite{schrimpf2019}. Dealers typically manage inventory risk through two complementary mechanisms. First, they engage in \emph{internalization}: skewing their quoted prices to attract flow that offsets existing positions and to discourage flow that would increase risk. Second, they perform \emph{externalization}: hedging residual risk by trading on the dealer-to-dealer (D2D) segment of the market or on all-to-all platforms \cite{butz2019}. The balance between these two mechanisms is a key strategic decision, with top-tier banks reportedly internalizing roughly 80\% of their flow in the so-called G10 currencies, the ten most heavily traded currencies, comprising USD, EUR, JPY, GBP, AUD, CAD, CHF, NZD, NOK, and SEK, according to \cite{schrimpf2019}.

The academic study of optimal market making has a rich history. The foundational work of Ho and Stoll \cite{ho1981} established the basic framework for dealer pricing under inventory and return uncertainty. This line of research was reinvigorated for modern electronic markets by Avellaneda and Stoikov \cite{avellaneda2008}, who formulated the market making problem as a stochastic optimal control problem in the context of a limit order book. Their approach leads to a Hamilton-Jacobi-Bellman (HJB) equation whose solution determines the optimal bid and ask quotes as functions of the dealer's inventory and time. Gu\'{e}ant, Lehalle, and Fernandez-Tapia \cite{gueant2013} subsequently provided a more detailed analysis of this control problem, including closed-form approximations to the optimal quotes. In parallel, Cartea, Jaimungal, and Ricci \cite{cartea2014} proposed a mean-variance-style objective that offers a more intuitive interpretation of the risk-reward trade-off, and this framework was further developed in the comprehensive treatments by Cartea, Jaimungal, and Penalva \cite{cartea2015} and Gu\'{e}ant \cite{gueant2016, gueant2017}.

A significant challenge arises when extending these models from a single asset to a portfolio of assets. In practice, market makers quote dozens or hundreds of instruments simultaneously, and the correlations between them are central to risk management. The multi-asset market making problem was addressed by Gu\'{e}ant \cite{gueant2016, gueant2017}, but computing the optimal strategy in general requires solving a high-dimensional HJB equation, where the dimensionality grows with the number of assets. This is a manifestation of the well-known \emph{curse of dimensionality}: grid-based numerical methods become computationally infeasible as the state space dimension increases. To address this, Bergault, Evangelista, Gu\'{e}ant, and Vieira \cite{bergault2021} proposed a technique that reduces the problem to a
linear-quadratic setting by approximating the Hamiltonian functions with quadratic forms. This yields a system of matrix Riccati ordinary differential equations (ODEs) whose dimension grows only quadratically in the number of assets, making the approach tractable even for large portfolios.

However, the FX market has a distinctive structure that sets it apart from other asset classes. In equity or bond markets, all assets are denominated in a single currency, and inventory is naturally measured in that currency. In FX, each traded instrument is an exchange rate between two currencies, and inventory management must be understood at the level of individual currency exposures. Moreover, the presence of \emph{cross pairs}, currency pairs that do not involve the reference currency, introduces both redundancy and opportunity. For instance, buying EUR/GBP is economically equivalent to simultaneously buying EUR/USD and selling GBP/USD. This triangular structure means that the dealer's positions in different currency pairs are inherently linked, and risk management must account for these interdependencies.

The model of Barzykin, Bergault, and Gu\'{e}ant \cite{barzykin2023} is, to our knowledge, the first to address this multi-currency structure in a comprehensive market making framework. By working with currency inventories valued in a reference currency (typically USD) as the relevant state variables, they formulate the dealer's optimization problem as a stochastic control problem that naturally accounts for cross-pair effects, correlations between exchange rates, client tiering, pricing ladders at multiple trade sizes, and external hedging with transaction costs and market impact. The resulting HJB equation is high-dimensional, but the authors adapt the quadratic approximation technique of \cite{bergault2021} to obtain a system of matrix Riccati-like ODEs whose dimension scales with the number of \emph{currencies} rather than the number of \emph{currency pairs}. This makes the approach computationally feasible even for markets with many currencies. The authors validate their approximation against grid-based PDE solutions in the two-currency case and report good agreement except under extreme and practically less relevant parameter regimes.

While the ODE-based approximation of \cite{barzykin2023} is elegant and scalable, several questions remain open. The validation against the full HJB solution is presented only briefly, and a systematic study of when and why the approximation breaks down has not been carried out. More broadly, the sensitivity of the optimal market making strategy to changes in model parameters, such as risk aversion, exchange rate volatilities, cross-pair correlations, and order flow intensities, has not been thoroughly explored. Understanding these sensitivities is important for at least two reasons. First, from a practical standpoint, the parameters entering the model are estimated from market data and are therefore subject to uncertainty; a strategy that is highly sensitive to a poorly estimated parameter may be unreliable in practice. This concern connects to the broader literature on robustness and parameter ambiguity in market making, as studied by Cartea, Donnelly, and Jaimungal \cite{cartea2017} in the context of robust market making under model uncertainty. Second, from a modelling perspective, sensitivity analysis reveals which features of the market environment are most consequential for the dealer's behavior, deepening our understanding of the economic mechanisms at play.

\section*{Contributions and outline}

In this thesis, we develop a grid-based numerical solver for the HJB equation arising in the multi-currency market making model of \cite{barzykin2023}, focusing on low-dimensional settings involving two to three currencies. In the two-currency case, we use the numerically computed PDE solution as a benchmark to evaluate the accuracy of the ODE-based approximation proposed in \cite{barzykin2023}. This serves the dual purpose of validating our numerical solver against a known reference solution and providing a more detailed assessment of the approximation quality across a range of parameter regimes.

Building on these numerical solutions, the main contribution of this thesis is a systematic sensitivity analysis and uncertainty quantification study. We investigate how key outputs of the model, including optimal bid-ask spreads, skew profiles, internalization zones, and hedging rates, respond to changes in the model parameters. In particular, we study the roles of risk aversion, exchange rate volatility, cross-pair correlations, and the characteristics of the client order flow. By doing so, we aim to identify which parameters most strongly influence the optimal strategy, to characterize the regimes in which the ODE approximation remains reliable, and to provide practical guidance on the robustness of the model's prescriptions under parameter uncertainty.

The remainder of this thesis is organized as follows. Chapter~\ref{ch:model} presents the multi-currency market making model of \cite{barzykin2023} in detail, including the stochastic control formulation and the derivation of the HJB equation. Chapter~\ref{ch:ode} derives the ODE-based approximation via a quadratic ansatz, describes the resulting Riccati system and its solution, and reproduces the numerical results of \cite{barzykin2023}. Chapter~\ref{ch:sensitivity} presents a systematic sensitivity analysis and uncertainty quantification study on the ODE model, using Sobol indices to identify the most influential parameters and forward uncertainty propagation to characterize the output distributions. Chapter~\ref{ch:pde} develops a grid-based PDE solver for the full HJB equation using fully implicit Euler time stepping with policy iteration, and compares the PDE and ODE solutions across the parameter regimes identified in the sensitivity analysis. Finally, Chapter~\ref{ch:conclusion} summarizes the findings and discusses directions for future work.